% 第 3 章
\bichapter{布局与优化}{Placement and Optimization}
\label{chap:place}

要了解如何在 IC Compiler II 中执行布局与优化，请参阅：
\begin{itemize}
  \item 执行布局与优化
  \item 识别优化期间无法修复的问题
  \item 分析布局
  \item 分析时序
  \item 分析功耗
  \item 对比 QoR 数据
\end{itemize}

% ============================================================
\section{执行布局与优化}
\subsection*{Performing Placement and Optimization}

完成设计规划与电源规划后，即可对设计执行布局（placement）、优化（optimization）与合法化（legalization）。以下主题说明如何进行：
\begin{itemize}
  \item 执行独立布局与合法化
  \item 使用 \cmd{place\_opt} 命令执行布局与优化
  \item 使用 Design Compiler 工具的物理引导（physical guidance）
  \item 执行多比特寄存器优化
  \item 执行 Magnet Placement
  \item 细化布局（Refining Placement）
  \item 对多电压 block 执行布局与优化
  \item 重建缓冲树（Buffer Trees）
  \item 添加与移除 tie cell
\end{itemize}

% ------------------------------------------------------------
\subsection{执行独立布局与合法化}
\subsubsection*{Performing Standalone Placement and Legalization}

要执行粗布局（coarse placement），使用 \cmd{create\_placement} 命令。若 block 含有通过读取 SCANDEF 文件标注的 scan chain，粗布局器还会执行 scan chain 优化。

可通过 \opt{-effort} 选项控制运行时间与 QoR 之间的权衡。

要启用：
\begin{itemize}
  \item 时序驱动布局：使用 \opt{-timing\_driven} 选项
  \item 考虑高扇出 net 与长 net，并计入优化阶段随后在这些 net 上添加的 buffer 的时序驱动布局：使用 \opt{-buffering\_aware\_timing\_driven} 选项
  \item 拥塞驱动布局：使用 \opt{-congestion} 选项\\
        要控制用于拥塞消除的 CPU 努力程度，使用 \opt{-congestion\_effort} 选项
  \item 拥塞驱动重构：使用 \opt{-congestion\_driven\_restructuring} 选项\\
        使用该选项时，工具会进行多轮布局与逻辑重构，以帮助降低拥塞
\end{itemize}

\begin{noteBox}
在使用 \cmd{create\_placement} 运行布局之前，应先运行 \cmd{merge\_clock\_gates} 合并 clock-gating 逻辑。
\end{noteBox}

要执行合法化，使用 \cmd{legalize\_placement} 命令。有关可指定的合法化约束与设置的更多信息，见「指定合法化设置」。

\begin{noteBox}
按「启用高级合法化算法」所述启用 advanced legalizer 时，不支持 \cmd{legalize\_placement} 的 \opt{-cells} 选项。
\end{noteBox}

% ------------------------------------------------------------
\subsection{使用 place\_opt 命令执行布局与优化}
\subsubsection*{Performing Placement and Optimization With the place\_opt Command}

要用单一命令执行粗布局、物理优化与合法化，使用 \cmd{place\_opt} 命令。

该命令支持多线程，并使用由 \cmd{set\_host\_options -max\_cores} 指定的线程数。

\cmd{place\_opt} 命令包含以下阶段：
\begin{enumerate}
  \item \textbf{初始布局（\cmd{initial\_place}）}\\
        此阶段工具合并 clock-gating 逻辑并执行粗布局。若 block 含有通过读取 SCANDEF 标注的 scan chain，工具还会执行 scan chain 优化。
  \item \textbf{初始 DRC 违例修复（\cmd{initial\_drc}）}\\
        此阶段工具移除已有缓冲树，并执行高扇出 net 综合（high-fanout-net synthesis）与电气 DRC 违例修复。
  \item \textbf{初始优化（\cmd{initial\_opto}）}\\
        此阶段工具执行时序、面积、拥塞与漏电功耗优化。
  \item \textbf{最终布局（\cmd{final\_place}）}\\
        此阶段工具执行增量布局以改善时序与拥塞，并对设计合法化。
  \item \textbf{最终优化（\cmd{final\_opto}）}\\
        此阶段工具进一步优化并合法化，以改善时序与拥塞。
\end{enumerate}

运行 \cmd{place\_opt} 时，默认执行全部布局与优化阶段。

要将 \cmd{place\_opt} 限制为仅运行其中一个或多个阶段，使用 \opt{-from} 指定起始阶段，使用 \opt{-to} 指定结束阶段（之后停止）。

若不指定 \opt{-from}，工具从 \cmd{initial\_place} 阶段开始；若不指定 \opt{-to}，工具一直运行到 \cmd{final\_opto} 阶段完成。

例如，要跳过初始布局，但对其余优化与布局阶段在已完成初始布局的设计上运行，可使用：
\begin{lstlisting}
icc2_shell> place_opt -from initial_drc
\end{lstlisting}

\subsubsection{为布局与优化创建临时时钟树}
\subsubsection*{Creating a Temporary Clock Tree for Placement and Optimization}

默认情况下，\cmd{place\_opt} 在布局与优化期间使用理想时钟（ideal clock）信息。但可通过将应用选项 \appopt{place\_opt.flow.trial\_clock\_tree} 设为 \cmd{true}，让工具构建临时时钟树，并在 \cmd{place\_opt} 期间使用传播时钟（propagated clocks）。

布局与优化之后，使用 \cmd{clock\_opt} 或 \cmd{synthesize\_clock\_trees} 执行时钟树综合时，工具会移除临时时钟树并构建实际时钟树。为使临时时钟树与实际时钟树相关，应在运行 \cmd{place\_opt} 之前指定时钟树综合设置。

\subsubsection{优化 Clock-Gating 单元}
\subsubsection*{Optimizing Clock-Gating Cells}

默认情况下，工具在以下时机合并 clock-gating 单元：
\begin{itemize}
  \item 未启用 Synopsys physical guidance 流程时，在 \cmd{place\_opt} 的 \cmd{initial\_place} 阶段开始处；
  \item 通过将 \appopt{place\_opt.flow.do\_spg} 设为 \cmd{true} 启用 Synopsys physical guidance（SPG）流程时，在 \cmd{place\_opt} 的 \cmd{initial\_opt} 阶段开始处。
\end{itemize}

但若在运行 \cmd{place\_opt} 之前已运行 \cmd{merge\_clock\_gates}，则 \cmd{place\_opt} 期间不再合并 clock gate。

可按如下方式控制 clock-gating 单元的合并：
\begin{itemize}
  \item 仅合并处于时钟树同一层级的 clock-gating 单元：将 \appopt{cts.icg.merge\_cross\_level} 设为 \cmd{false}。\\
        默认情况下，工具可合并处于时钟树不同层级的 clock-gating 单元。
  \item 禁用 clock-gating 单元合并：将 \appopt{place\_opt.flow.merge\_clock\_gates} 设为 \cmd{false}。
\end{itemize}

若 clock gate 的 enable 引脚处于关键时序路径上，\cmd{place\_opt} 还可进一步优化 clock-gating 单元。要启用该功能，将 \appopt{place\_opt.flow.optimize\_icgs} 设为 \cmd{true}。

在 clock-gating 单元优化期间，工具会：
\begin{itemize}
  \item 执行 clock-aware 布局（无论 \appopt{place\_opt.flow.clock\_aware\_placement} 如何设置），并将关键 clock-gating 单元及其扇出放置在更优位置；
  \item 构建临时时钟树（无论 \appopt{place\_opt.flow.trial\_clock\_tree} 如何设置），并用该时钟树识别时序关键的 clock-gating 单元；
  \item 若 clock gate 的 enable 引脚处于关键时序路径上，则拆分（split）clock-gating 单元。\\
        可通过 \appopt{place\_opt.flow.optimize\_icgs\_critical\_range} 指定临界范围，以控制拆分程度。
\end{itemize}

\subsubsection{对 Clock Gate 使用更准确的 Latency}
\subsubsection*{Using Accurate Latencies for Clock Gates}

默认情况下，工具在整个 \cmd{place\_opt} 流程中估计并更新 clock-gate latency。这确保在时钟树综合之前、时钟仍为理想时钟时，数据通路优化与 useful-skew 优化能使用更准确的 clock-gate latency。

要提高估计 clock-gate latency 的准确性，应在运行 \cmd{place\_opt} 之前指定时钟树综合设置，例如时钟树 buffer/inverter、时钟非默认布线规则（nondefault routing rules）等。

\begin{noteBox}
若已按「为布局与优化创建临时时钟树」启用 trial clock tree synthesis，或按「优化 Clock-Gating 单元」启用 clock-gate 优化（二者在 CTS 之前用更准确的方法建模时钟树效应），则工具不再执行 clock-gate latency 估计与 clock-gate-latency-aware 布局。
\end{noteBox}

工具不会修改你对 clock-gating 单元指定的任何 \cmd{set\_clock\_latency} 约束。工具估计的 clock-gate latency 值通过 \cmd{set\_clock\_latency -offset} 存储。可用 \cmd{write\_script -format icc2} 查看这些估计值，并在生成输出中查找 \cmd{set\_clock\_latency -offset} 命令。\cmd{write\_sdc} 生成的 SDC 文件不会捕获这些估计的 latency offset。

要阻止工具估计特定 clock-gate 的 latency，对其时钟引脚设置 \cmd{dont\_estimate\_clock\_latency} 属性，例如：
\begin{lstlisting}
icc2_shell> set_attribute \
   [get_pins {ICG21/CK}] dont_estimate_clock_latency true
\end{lstlisting}

\subsubsection{启用基于全局布线的高扇出综合}
\subsubsection*{Enabling Global Route Based High-Fanout Synthesis}

工具执行高扇出综合与电气 DRC 违例修复时，默认使用虚拟布线（virtual routes）。

要在高扇出综合与电气 DRC 违例修复期间使用全局布线（global routes），在运行 \cmd{place\_opt} 之前设置 \appopt{place\_opt.initial\_drc.global\_route\_based}：
\begin{lstlisting}
icc2_shell> set_app_options \
   -name place_opt.initial_drc.global_route_based -value 1
\end{lstlisting}

要启用包含 trial 高扇出综合的集成两遍初始布局流程，在运行 \cmd{place\_opt} 之前设置：
\begin{lstlisting}
icc2_shell> set_app_options \
   -name place_opt.initial_drc.global_route_based -value 1
icc2_shell> set_app_options \
   -name place_opt.initial_place.two_pass -value true
\end{lstlisting}

该集成两遍初始布局流程包括：线长驱动布局、trial 高扇出综合，以及时序驱动布局。工具丢弃 trial 高扇出综合结果，但保留时序驱动布局，并继续进行基于全局布线的高扇出综合。

对 floorplan 碎片化严重的设计，启用集成两遍初始布局与基于全局布线的高扇出综合，有助于改善时序并降低拥塞。

\subsubsection{启用缓冲树重建}
\subsubsection*{Enabling the Rebuilding of Buffer Trees}

在 \cmd{place\_opt} 的 \cmd{initial\_drc} 阶段，工具默认执行高扇出 net 综合，为高扇出 net 构建缓冲树。\cmd{initial\_drc} 之后，若能改善 setup 时序、逻辑 DRC 违例、线长或面积，工具可移除并重建整棵缓冲树或部分缓冲树。

要在以下阶段启用缓冲树重建：
\begin{itemize}
  \item 在 \cmd{initial\_opto} 阶段针对多电压 net：将 \appopt{opt.buffering.enable\_mv\_rebuffering} 设为 \cmd{true}
  \item 在 \cmd{final\_opto} 阶段针对任意 net：将 \appopt{opt.buffering.enable\_rebuffering} 设为 \cmd{true}
\end{itemize}

\subsubsection{更改拥塞努力程度}
\subsubsection*{Changing the Congestion Effort}

默认情况下，\cmd{place\_opt} 执行拥塞优化以改善设计的可布线性。可通过 \appopt{place\_opt.place.congestion\_effort} 禁用或控制该优化的努力程度；默认值为 \cmd{medium}。

\subsubsection{对关键 Net 使用非默认布线规则}
\subsubsection*{Using of Nondefault Routing Rules for Critical Nets}

为改善时序 QoR，工具可在布线前优化期间对时序关键 net 使用非默认布线规则（nondefault routing rules）。此外，工具可引导路由器将这些非默认规则分配作为软约束加以遵守。

要为 \cmd{place\_opt} 启用该能力，将 \appopt{place\_opt.flow.optimize\_ndr} 设为 \cmd{true}：
\begin{lstlisting}
icc2_shell> set_app_options -name place_opt.flow.optimize_ndr \
    -value true
\end{lstlisting}

\subsubsection{启用路径优化}
\subsubsection*{Enabling Path Optimization}

路径优化（path optimization）是一种增量优化能力，专注于改善时序路径的 QoR。也可在 \cmd{place\_opt} 期间通过将 \appopt{place\_opt.flow.do\_path\_opt} 设为 \cmd{true} 来启用：
\begin{lstlisting}
icc2_shell> set_app_options -name place_opt.flow.do_path_opt -value true
\end{lstlisting}

\subsubsection{在 place\_opt 期间执行 IR-Drop-Aware 布局}
\subsubsection*{Performing IR-Drop-Aware Placement During the place\_opt Command}

布局期间，工具可利用单元的电压（IR）drop 值识别高功耗密度区域，并分散高压降单元，从而降低此类区域的功耗密度。

要在 \cmd{place\_opt} 期间执行 IR-drop-aware 布局，按下列步骤：
\begin{enumerate}
  \item 用 \cmd{place\_opt -to final\_place} 运行到最终布局阶段。
  \item 为 RedHawk Fusion 做好设置，并用 \cmd{analyze\_rail -voltage\_drop} 执行静态或动态压降分析，例如：
\begin{lstlisting}
icc2_shell> source redhawk_setup.tcl
icc2_shell> analyze_rail -voltage_drop static -nets {VDD VSS}
\end{lstlisting}
        更多信息见「执行压降分析」。
  \item 将 \appopt{place.coarse.ir\_drop\_aware} 设为 \cmd{true}，以启用 IR-drop-aware 布局。
  \item （可选）按「控制 IR-Drop-Aware 布局」指定额外设置。
  \item 用 \cmd{place\_opt -from final\_place} 重新运行最终布局阶段。
\end{enumerate}

\subsubsection{在 place\_opt 期间执行并发时钟与数据优化}
\subsubsection*{Performing Concurrent Clock and Data Optimization During the place\_opt Command}

在数据通路优化中应用 useful skew 技术，利用正松弛并调整寄存器的时钟到达时间以改善时序 QoR，称为并发时钟与数据（concurrent clock and data，CCD）优化。

默认情况下，工具在 \cmd{place\_opt} 的最终优化阶段执行 CCD。可通过将 \appopt{place\_opt.flow.enable\_ccd} 设为 \cmd{false} 来禁用。

\cmd{place\_opt} 期间导出的时钟 latency 调整以 clock-latency offset 形式通过 \cmd{set\_clock\_latency -offset} 存储，并以 clock-balance-point-delay offset 形式通过 \cmd{set\_clock\_balance\_point -offset} 存储。可用 \cmd{write\_script -format icc2} 查看这些 clock-latency offset（在输出中查找 \cmd{set\_clock\_latency -offset}）。\cmd{write\_sdc} 生成的 SDC 不会捕获这些 offset。可用 \cmd{report\_clock\_balance\_point} 报告 clock-balance-point-delay offset。

要为 \cmd{place\_opt} 期间的 CCD 设置 block，执行下列步骤：
\begin{enumerate}
  \item 确保为 \cmd{clock\_opt} 设为 active 的 scenario，对 \cmd{place\_opt} 也保持 active。
  \item 通过下列两种设置之一，指定工具在 \cmd{place\_opt} 期间构建临时时钟树并使用传播时钟：
\begin{lstlisting}
icc2_shell> set_app_options -name place_opt.flow.trial_clock_tree \
     -value true
\end{lstlisting}
\begin{lstlisting}
icc2_shell> set_app_options -name place_opt.flow.optimize_icgs \
     -value true
\end{lstlisting}
        或者，工具可用理想时钟做 useful skew 计算；但必须为所有 clock-gating 单元与时序单元（含 macro）的时钟引脚指定准确的理想时钟 latency。
\end{enumerate}

可通过下列可选步骤改变 \cmd{place\_opt} 期间 CCD 的默认行为：
\begin{enumerate}
  \item （可选）限制 CCD 的 latency 调整值，见「限制 Latency 调整值」。
  \item （可选）控制边界寄存器的 latency 调整，见「排除边界路径」。
  \item （可选）在 CCD 中忽略特定 path group，见「排除特定 Path Group」。
  \item （可选）在 CCD 中忽略特定 scenario，见「排除特定 Scenario」。
  \item （可选）在 CCD 中忽略特定 sink，见「排除特定 Sink」。
  \item （可选）分析 CCD 效果，见「报告并发时钟与数据时序」。
  \item （可选）对带内部延时的 macro 导出自动零平衡点（automatic zero balance point）。
  \item （可选）将 latch 与 discrete clock gate 按相同 offset 一起 skew，见「Skewing Latch and Discrete Clock Gates」。
\end{enumerate}

% ------------------------------------------------------------
\subsection{使用 Design Compiler 工具的物理引导}
\subsubsection*{Using Physical Guidance From the Design Compiler Tool}

使用 \cmd{place\_opt} 执行布局与优化时，可将 Design Compiler Graphical 工具的 Synopsys physical guidance 信息作为起点。

使用 Design Compiler Graphical 布局的好处包括：
\begin{itemize}
  \item 缩短布局步骤的运行时间；
  \item 改善 Design Compiler Graphical 与 IC Compiler II 之间的相关性。
\end{itemize}

要从 Design Compiler Graphical 传递设计，使用 Verilog 格式的门级网表以及约束信息。

为保持一致性并改善相关性，必须：
\begin{itemize}
  \item 在两个工具中使用相同的逻辑约束（含时序与 UPF 信息）。\\
        但须将 Design Compiler 的 scenario 约束转换为 IC Compiler II 支持的 mode、corner 与 scenario 格式。
  \item 在 IC Compiler II 中使用 Design Compiler Graphical 布局，并在两个工具中使用相同的 floorplan 与物理约束。\\
        不过，可提供 Design Compiler Graphical 当前不支持、但 IC Compiler II 物理实现所需的额外 floorplan 信息。
\end{itemize}

要从 Design Compiler Graphical 传递物理信息（布局、floorplan 与物理约束），执行下列任务：
\begin{itemize}
  \item 在 Design Compiler Graphical 中用 \cmd{write\_def} 生成 DEF 文件，并在 IC Compiler II 中用
\begin{lstlisting}
read_def -add_def_only_objects {cells} -convert_sites
\end{lstlisting}
        读入。\\
        该 DEF 包含：
        \begin{itemize}
          \item Die area、site rows 与 tracks
          \item Port 的位置与形状
          \item Macro、标准单元与 physical-only 单元的位置
          \item Route guide 与预布线 net
          \item Placement blockage
        \end{itemize}
  \item 创建 Tcl 文件，包含 DEF 未捕获的物理约束，并在 IC Compiler II 中应用到设计。\\
        DEF 未捕获的物理约束包括：
        \begin{itemize}
          \item Voltage area 定义
          \item Layer 约束
          \item 带利用率的特殊 route guide
          \item 带 blocked layers 的特殊 blockage
        \end{itemize}
\end{itemize}

要在 IC Compiler II 中使用 Synopsys physical guidance，在运行 \cmd{place\_opt} 之前设置：
\begin{lstlisting}
icc2_shell> set_app_options -name place_opt.flow.do_spg -value true
\end{lstlisting}

下列示例脚本展示如何将 Synopsys physical guidance 与 \cmd{place\_opt} 配合使用：
\begin{lstlisting}
# Read in the netlist
read_verilog design_spg.v

# Apply the timing constraints
source timing_const.tcl

# Apply the UPF constraints
load_upf design.upf

# Apply the DEF file generated from the Design Compiler Graphical tool
read_def -add_def_only_objects {cells} -convert_sites design_dcg.def

# Apply the placement constraints not supported by DEF
source placement_const.tcl

# Enable the Synopsys placement guidance flow
set_app_options -name place_opt.flow.do_spg -value true

# Perform placement and optimization
place_opt
\end{lstlisting}

更多信息见 \textit{Design Compiler User Guide} 中的 Synopsys physical guidance。

% ------------------------------------------------------------
\subsection{执行多比特寄存器优化}
\subsubsection*{Performing Multibit Register Optimization}

IC Compiler II 可将单比特寄存器或较小的多比特寄存器组合（bank），并替换为等效的更大多比特寄存器。例如，可将八个 1-bit 寄存器或四个 2-bit 寄存器 bank 组合，替换为一个 8-bit 寄存器 bank 或两个 4-bit 寄存器 bank。工具仅在单比特寄存器具有相同时序约束时才合并它们，并将时序约束复制到结果多比特寄存器。

\begin{figure}[htbp]
\centering
\fbox{\parbox{0.85\textwidth}{\centering\vspace{1.5cm}
\textit{（原书 Figure 25：Replacing Multiple Single-Bit Register Cells With a Multibit Register Cell）}\\[0.3em]
\small 请参阅原版 PDF 插图；可将截图放入 \texttt{figures/} 目录后用 \texttt{\textbackslash includegraphics} 替换。
\vspace{1.5cm}}}
\caption{用多比特寄存器单元替换多个单比特寄存器单元}
\label{fig:multibit-replace}
\end{figure}

用多比特寄存器替换单比特寄存器可减少：
\begin{itemize}
  \item 面积（因共享晶体管与优化的晶体管级版图）；
  \item 时钟树总网线长度；
  \item 时钟树 buffer 数量与时钟树功耗。
\end{itemize}

若能改善总负松弛（total negative slack），工具也可将大多比特寄存器拆分（debank）为较小的多比特寄存器或单比特寄存器。

某些逻辑库具有混合驱动强度（mixed-drive-strength）的多比特寄存器，部分 bit 驱动强度高于其他 bit。对此类单元，若违例路径经过较低驱动强度的 bit，工具可对混合驱动强度多比特单元重新布线（rewire），使违例路径经过较高驱动强度的 bit。

以下主题说明如何执行多比特寄存器优化：
\begin{itemize}
  \item 执行集成多比特寄存器优化
  \item 使用离散命令执行多比特寄存器优化
  \item 对多比特保持寄存器（retention registers）做 banking
\end{itemize}

\subsubsection{执行集成多比特寄存器优化}
\subsubsection*{Performing Integrated Multibit Register Optimization}

要在 \cmd{place\_opt} 期间执行多比特寄存器优化：
\begin{enumerate}
  \item 将 \appopt{place\_opt.flow.enable\_multibit\_banking} 设为 \cmd{true}，启用多比特寄存器 banking。
  \item （可选）将 \appopt{place\_opt.flow.enable\_multibit\_debanking} 设为 \cmd{true}，启用 debanking。
  \item （可选）将 \appopt{multibit.banking.across\_equivalent\_icg} 设为 \cmd{true}，允许对由等效 integrated-clock-gating（ICG）单元驱动的寄存器做 banking。\\
        默认情况下，工具仅对由同一时钟 net 驱动的寄存器做 banking。启用后，工具会对等效 clock-gating 单元驱动的寄存器做 banking，并合并这些 clock-gating 单元。但这样做可能丢失施加在连接至 clock-gating 单元的 net 上的非默认布线规则。因此优化后请检查并重新应用这些规则。
  \item （可选）将 \appopt{place\_opt.flow.enable\_multibit\_rewiring} 设为 \cmd{true}，启用多比特寄存器 rewiring。
  \item （可选）用 \cmd{set\_multibit\_options} 指定 banking 设置。例如，将 \cmd{reg[1]} 排除在多比特优化之外：
\begin{lstlisting}
icc2_shell> set_multibit_options -exclude [get_cells "reg[1]"]
\end{lstlisting}
  \item （可选）按表~\ref{tab:multibit-appopt} 控制 banking 与 debanking。
  \item 用 \cmd{place\_opt} 执行布局与优化。
  \item 用 \cmd{report\_multibit} 报告多比特组件。
\end{enumerate}

\begin{table}[htbp]
\centering
\caption{控制多比特寄存器 Banking 与 Debanking 的应用选项}
\label{tab:multibit-appopt}
\begin{tabularx}{\textwidth}{@{}X X@{}}
\toprule
\textbf{目的} & \textbf{将下列应用选项设为 true} \\
\midrule
Banking 时忽略未包含在 SCANDEF 信息中的 scan flip-flop
  & \appopt{multibit.banking.enable\_strict\_scan\_check} \\
Banking 时考虑带多个 scan-in 引脚与内部 scan-in 引脚的多比特寄存器库单元
  & \appopt{multibit.banking.lcs\_consider\_multi\_si\_cells} \\
Banking 与 debanking 时忽略用户指定的 size-only 单元
  & \appopt{multibit.common.ignore\_sizeonly\_cells} \\
\bottomrule
\end{tabularx}
\end{table}

\subsubsection{使用离散命令执行多比特寄存器优化}
\subsubsection*{Performing Multibit Register Optimization Using Discrete Commands}

可用下列离散步骤执行多比特 banking：
\begin{enumerate}
  \item 用 \cmd{create\_placement} 执行初始布局。
  \item （可选）用 \cmd{set\_multibit\_options} 指定 banking 设置。
  \item 用 \cmd{identify\_multibit} 识别可替换为多比特单元的单元组，见「识别多比特 Bank」。
  \item 用 \cmd{place\_opt -from initial\_drc} 优化设计。
  \item （可选）用 \cmd{report\_multibit} 报告 banking 信息。
  \item （可选）若 banking 未改善 QoR，用 \cmd{split\_multibit} 将特定多比特 bank 拆分为单独寄存器或更小的多比特 bank，见「拆分多比特 Bank」。
\end{enumerate}

\paragraph{识别多比特 Bank}
\paragraph*{Identifying Multibit Banks}

要用 \cmd{identify\_multibit} 识别可替换为多比特单元的寄存器、isolation cell 或 level-shifter 组，必须指定：
\begin{itemize}
  \item 如何应用 banking 信息，任选其一：
        \begin{itemize}
          \item 用 \opt{-apply} 让工具将 banking 信息应用到设计；
          \item 用 \opt{-output\_file} 生成包含 \cmd{create\_multibit} 命令的输出文件，以便手动应用。\\
                可查看、必要时编辑该输出，再用 \cmd{source} 应用到设计。
        \end{itemize}
  \item 要 bank 的单元类型，任选其一：
        \begin{itemize}
          \item 用 \opt{-register} 识别可 bank 的寄存器组；
          \item 用 \opt{-mv\_cell} 识别可 bank 的 isolation 或 level-shifter 组。
        \end{itemize}
\end{itemize}

使用 \opt{-register} banking 寄存器时，可选地：
\begin{itemize}
  \item 用 \opt{-input\_map\_file} 指定输入映射文件，控制单比特到多比特 bank 的映射。该文件包含哪些多比特寄存器 bank 替换哪些单比特寄存器，例如：
\begin{lstlisting}
reg_group_1 {REGX1 REGX2 REGX4}
2 {1 MREG2}
4 {1 MREG4}
6 {1 MREG2}{1 MREG4}
\end{lstlisting}
        该映射文件包含名为 \cmd{reg\_group\_1} 的功能组，说明如何将参考单元类型为 REGX1、REGX2、REGX4 的单比特寄存器组合为多比特寄存器：
        \begin{itemize}
          \item 第一行：两个单比特寄存器可组合成一个由一个 MREG2 参考单元构成的寄存器 bank；
          \item 第二行：四个单比特寄存器可组合成一个由一个 MREG4 参考单元构成的寄存器 bank；
          \item 第三行：六个单比特寄存器可组合成由一个 MREG2 与一个 MREG4 构成的寄存器 bank。
        \end{itemize}
  \item 用 \opt{-slack\_threshold} 按松弛排除寄存器。松弛小于指定阈值的寄存器被忽略。默认考虑所有寄存器。
\end{itemize}

此外，可选地通过下列方式控制被 bank 的单元：
\begin{itemize}
  \item 用 \opt{-cells} 指定要考虑的实例列表（使用后不考虑其他单元）；
  \item 用 \opt{-exclude\_instance} 指定要排除的实例列表；
  \item 用 \opt{-exclude\_library\_cells} 指定要排除其全部实例的库单元列表。
\end{itemize}

默认情况下，\cmd{identify\_multibit} 执行 DFT 优化，重新划分 scan chain，使由同一时钟驱动器或 gating 单元驱动的寄存器进入同一 scan chain，从而增加可组合成多比特寄存器 bank 的单比特寄存器数量。若不希望重新划分 scan chain，对 \cmd{identify\_multibit} 使用 \opt{-no\_dft\_opt}。

\paragraph{拆分多比特 Bank}
\paragraph*{Splitting Multibit Banks}

为改善局部拥塞或路径松弛，可用 \cmd{split\_multibit} 将多比特 bank 拆分为更小的多比特 bank 或单比特单元。默认情况下，该命令拆分负松弛时序路径上的所有多比特单元。

要拆分：
\begin{itemize}
  \item 特定单元：指定单元实例名。可用 \opt{-lib\_cells} 指定拆分后使用的较小单元的参考名。\\
        下例将名为 \cmd{reg\_gr0} 的四比特寄存器 bank 拆分为两个两比特多比特 bank：
\begin{lstlisting}
icc2_shell> split_multibit reg_gr0 -lib_cells {lib1/MREG2 lib1/MREG2}
\end{lstlisting}
  \item 松弛小于特定值的路径上的多比特单元：使用 \opt{-slack\_threshold}。\\
        下例拆分松弛差于 $-1$ 的路径上的多比特单元：
\begin{lstlisting}
icc2_shell> split_multibit -slack_threshold -1
\end{lstlisting}
  \item 特定时序 path group 中的多比特单元：使用 \opt{-path\_groups}。\\
        使用该选项时，工具拆分指定 path group 中负松弛路径上的多比特单元。可与 \opt{-slack\_threshold} 联用为该 path group 指定松弛阈值。\\
        下例拆分 path group CLKA 中松弛差于 $-0.5$ 的路径上的多比特单元：
\begin{lstlisting}
icc2_shell> split_multibit -path_groups CLKA \
   -slack_threshold -0.5
\end{lstlisting}
\end{itemize}

还可进一步控制拆分：
\begin{itemize}
  \item 用 \opt{-exclude\_instance} 排除特定单元；
  \item 用 \opt{-cells} 仅考虑特定单元。
\end{itemize}

\cmd{split\_multibit} 仅拆分由 \cmd{create\_multibit} 或 \cmd{identify\_multibit} 创建的多比特单元。

\subsubsection{对多比特保持寄存器做 Banking}
\subsubsection*{Banking Multibit Retention Registers}

工具可将单比特保持寄存器组合为多比特保持寄存器，若满足：
\begin{itemize}
  \item 单比特保持寄存器关联同一 retention strategy；
  \item 其 retention 引脚具有相同驱动器；
  \item 单元库中有匹配的多比特保持寄存器；
  \item 该多比特保持寄存器的库单元已通过 \cmd{map\_retention\_cell -lib\_cells} 指定为该 retention strategy 的库单元。
\end{itemize}

工具可拆分多比特保持寄存器，若：
\begin{itemize}
  \item 单元库中有匹配的单比特保持寄存器；
  \item 这些单比特保持寄存器的库单元已通过 \cmd{map\_retention\_cell -lib\_cells} 指定为该 retention strategy 要使用的库单元。
\end{itemize}

\begin{noteBox}
若使用 \cmd{identify\_multibit} 识别多比特保持寄存器，请使用 \opt{-register}，而非 \opt{-mv\_cell}。
\end{noteBox}

% ------------------------------------------------------------
\subsection{执行 Magnet Placement}
\subsubsection*{Performing Magnet Placement}

为改善复杂 floorplan 的拥塞或改善设计时序，可使用 magnet placement：将固定对象指定为磁体（magnet），并让工具将连接到该磁体对象的所有标准单元靠近其放置。磁体对象可为固定 macro 单元、固定 macro 的引脚，或 I/O port。

为获得最佳结果，应在标准单元布局之前执行 magnet placement。

执行 magnet placement 时，使用 \cmd{magnet\_placement} 命令，并指定磁体及所需特殊功能的选项。表~\ref{tab:magnet-opts} 列出 \cmd{magnet\_placement} 选项。

\begin{table}[htbp]
\centering
\caption{\cmd{magnet\_placement} 命令选项}
\label{tab:magnet-opts}
\begin{tabularx}{\textwidth}{@{}X l@{}}
\toprule
\textbf{目的} & \textbf{使用该选项} \\
\midrule
指定要拉向磁体对象的单元 & \opt{-cells cell\_list} \\
允许用 \opt{-cells} 指定层次单元，并说明如何处理这些层次单元
  & \opt{-hierarchy\_mode default\textbar block\textbar all} \\
允许移动标记为 fixed 的单元 & \opt{-move\_fixed} \\
允许移动标记为 legalize-only 的单元 & \opt{-move\_legalize\_only} \\
Magnet placement 后将移动过的单元标记为 fixed 或 legalize-only
  & \opt{-mark\_fixed} / \opt{-mark\_legalize\_only} \\
指定距磁体的逻辑级数，在该范围内考虑 magnet placement
  & \opt{-logical\_levels number} \\
确定逻辑级数时排除 buffer 与 inverter & \opt{-exclude\_buffers} \\
查找要拉取的单元时，在时序单元输入处停止路径追踪
  & \opt{-stop\_by\_sequential\_cells} \\
查找要拉取的单元时，在时序单元输出处停止路径追踪
  & \opt{-stop\_on\_sequential\_cells} \\
查找要拉取的单元时，在指定 pin、port 或单元之外停止路径追踪
  & \opt{-stop\_points} \\
防止将单元放在 soft blockage 上（默认仅避开 hard blockage）
  & \opt{-avoid\_soft\_blockage} \\
磁体为多个 port 或长 port 时，指定如何放置单元
  & \opt{-multiple\_long\_port\_mode auto\textbar nearby} \\
仅报告将被拉向磁体的单元，但不实际拉取
  & \opt{-only\_report\_magnet\_cells} \\
创建正在被拉向磁体的单元集合 & \opt{-get\_collection} \\
\bottomrule
\end{tabularx}
\end{table}

下列示例使用图~\ref{fig:magnet-datapath} 所示数据通路。

\begin{figure}[htbp]
\centering
\fbox{\parbox{0.85\textwidth}{\centering\vspace{1.2cm}
\textit{（原书 Figure 26：Example of Pulling Cells in a Contiguous Data Path）}
\vspace{1.2cm}}}
\caption{在连续数据通路中拉取单元的示例}
\label{fig:magnet-datapath}
\end{figure}

下列命令将所有单元 C1 至 C8 拉向磁体对象 C0：
\begin{lstlisting}
icc2_shell> magnet_placement C0
\end{lstlisting}

下列命令将 C6、C7、C8 拉向磁体对象 C2：
\begin{lstlisting}
icc2_shell> magnet_placement C2 -cells {C6 C7 C8}
\end{lstlisting}

尽管 C3、C4、C5 形成连续数据通路，但它们与 C0 隔离，因此下列命令不会将任何单元拉向 C0：
\begin{lstlisting}
icc2_shell> magnet_placement C0 -cells {C3 C4 C5}
\end{lstlisting}

下列命令也不会将任何单元拉向 C0，因为 C3、C6、C8 不形成连续数据通路：
\begin{lstlisting}
icc2_shell> magnet_placement C0 -cells {C3 C6 C8}
\end{lstlisting}

% ------------------------------------------------------------
\subsection{细化布局}
\subsubsection*{Refining Placement}

要细化布局并最小化拥塞，使用 \cmd{refine\_placement} 命令。默认对该命令对整个设计执行增量布局。
\begin{itemize}
  \item 要细化特定区域的布局，使用 \opt{-coordinates} 并指定区域坐标。
  \item 要更改布局努力程度，使用 \opt{-effort}，取值为 \cmd{low}、\cmd{medium} 或 \cmd{high}；默认为 \cmd{medium}。
  \item 要更改最小化拥塞的努力程度，使用 \opt{-congestion\_effort}，取值为 \cmd{low}、\cmd{medium} 或 \cmd{high}；默认为 \cmd{medium}。
  \item 要控制对现有布局的扰动程度，使用 \opt{-perturbation\_level}，取值为 \cmd{min}、\cmd{medium}、\cmd{high} 或 \cmd{max}；默认为 \cmd{medium}。
\end{itemize}

% ------------------------------------------------------------
\subsection{对多电压 Block 执行布局与优化}
\subsubsection*{Performing Placement and Optimization on Multivoltage Blocks}

布局与优化会自动考虑多电压约束，例如 voltage area。

布局期间，工具会：
\begin{itemize}
  \item 将单元放置在其关联的 voltage area 或 voltage area shape 内；
  \item 将 level shifter 与 isolation cell 靠近 voltage area 边界放置。
\end{itemize}

优化期间插入 buffer 时，工具会：
\begin{enumerate}
  \item 考虑 net 的驱动器与负载的相关供电，并用下列之一对电源状态做动态分析：
        \begin{itemize}
          \item 用 \cmd{create\_pst} 创建的 UPF power state table；
          \item 用 \cmd{add\_power\_state} 创建的 UPF power state groups。
        \end{itemize}
  \item 根据电源状态与被缓冲 voltage area 中的可用供电，选择不会引入新的多电压违例的缓冲供电。优先选用单轨（single-rail）buffer 而非双轨（dual-rail）buffer，以节省功耗与面积并改善可布线性。
\end{enumerate}

要了解特定 net 如何或是否被缓冲，使用 \cmd{check\_bufferability} 命令，见「分析 Net 的可缓冲性」。

% ------------------------------------------------------------
\subsection{重建缓冲树}
\subsubsection*{Rebuilding Buffer Trees}

在 \cmd{place\_opt} 的 \cmd{initial\_drc} 阶段，工具移除已有缓冲树并执行高扇出综合。运行 \cmd{place\_opt} 之后，要用 \cmd{create\_buffer\_trees} 按特定需求重建特定缓冲树。

默认情况下，\cmd{create\_buffer\_trees} 会：
\begin{itemize}
  \item 为所有高扇出 net 重建缓冲树。\\
        要重建特定高扇出 net，使用 \opt{-from}。
  \item 在重建前移除已有缓冲树。\\
        要保留已有 buffer 并增量构建缓冲树，使用 \opt{-incremental}。
  \item 使用虚拟布线拓扑与提取构建缓冲树。\\
        要基于全局布线构建缓冲树，使用 \opt{-global\_route\_based}。
  \item 不合法化新添加的单元。\\
        要合法化添加的单元，在 \cmd{create\_buffer\_trees} 之后运行 \cmd{legalize\_placement}。
\end{itemize}

下例仅为 \cmd{n33} net 重建缓冲树，且仅使用 \cmd{buf4x} 与 \cmd{buf8x} 类型 buffer：
\begin{lstlisting}
icc2_shell> set_lib_cell_purpose -include none {tech_lib/*}
icc2_shell> set_lib_cell_purpose \
   -include optimization {tech_lib/buf4x tech_lib/buf8x}
icc2_shell> create_buffer_trees -from n33 -global_route_based
\end{lstlisting}

\begin{noteBox}
\cmd{create\_buffer\_trees} 仅应在设计流程的布线前（preroute）阶段使用。
\end{noteBox}

% ------------------------------------------------------------
\subsection{添加与移除 Tie Cell}
\subsubsection*{Adding and Removing Tie Cells}

运行 \cmd{place\_opt} 或 \cmd{clock\_opt} 时，工具可在优化期间添加 tie cell。此外，可用 \cmd{add\_tie\_cells} 为特定 pin 或 port 添加 tie cell。\cmd{place\_opt}、\cmd{clock\_opt} 或 \cmd{add\_tie\_cells} 首次向 block 添加 tie cell 时，工具会优化 tie cell 连接：移除所有已有 tie cell，并添加经最优聚类的 tie cell。

要移除特定 tie cell，使用 \cmd{remove\_tie\_cells}。默认移除当前 block 中的所有 tie cell。要移除特定实例或某库单元的全部实例，使用 \opt{-objects}。

% ============================================================
\section{识别优化期间无法修复的问题}
\subsection*{Identifying Issues That Cannot Be Fixed During Optimization}

IC Compiler II 提供下列能力，用于识别优化期间无法修复的问题：
\begin{itemize}
  \item 分析 Net 的可缓冲性
  \item 分析无法修复的违例
\end{itemize}

\subsection{分析 Net 的可缓冲性}
\subsubsection*{Analyzing the Bufferability of Nets}

要分析当前 block 中的 net 是否可被缓冲，或报告为何无法缓冲，使用 \cmd{check\_bufferability}。该命令检查是否存在：
\begin{itemize}
  \item 在指定 voltage area 中可用于缓冲该 net 的 supply net；
  \item 可用于缓冲该 net 的合适库单元；
  \item 阻止缓冲的 net 设置，例如 ideal-network 或 don’t-touch 约束。
\end{itemize}

使用该命令时必须指定：
\begin{itemize}
  \item 要分析的 net，任选其一：
        \begin{itemize}
          \item 用 \opt{-nets} 指定 net 名；
          \item 用 \opt{-driver} 与 \opt{-loads} 指定驱动器与负载，并用 \opt{-hierarchy} 指定进行分析的逻辑层次。
        \end{itemize}
  \item 用 \opt{-voltage\_area} 指定分析可缓冲性的 voltage area。
\end{itemize}

例如，要分析名为 \cmd{n34} 的 net 在 PD1 电源域中是否可添加 buffer：
\begin{lstlisting}
icc2_shell> check_bufferability -nets n34 -voltage_area PD1
Information: Can insert dual rail buffers (8 lib_cells) and inverters (6
 lib_cells) with supply nets (power: VDDSW, ground: VSS). (MV-454)
\end{lstlisting}

要分析在 PD\_TOP voltage area 的 top 层次中，是否可在驱动引脚 \cmd{top\_reg\_1\_/Q} 与负载引脚 \cmd{top\_U1\_9/B} 之间添加 buffer：
\begin{lstlisting}
icc2_shell> check_bufferability -driver top_reg_1_/Q \
   -loads top_U1_9/B -hierarchy / -voltage_area PD_TOP
Information: Cannot insert buffers or inverters, because required supply
 nets are not available. (MV-459)
\end{lstlisting}

要获取关于 net 可缓冲性的更多信息，使用 \opt{-verbose}。

\subsection{分析无法修复的违例}
\subsubsection*{Analyzing Violations That Cannot Be Fixed}

要分析 block 并生成包含无法修复违例信息的报告，使用 \cmd{analyze\_design\_violations}。报告先给出分析摘要，再按降序列出违例明细。违例按无法修复的原因分类。

使用该命令时，必须用 \opt{-type} 指定要分析的违例类型：
\begin{itemize}
  \item \cmd{max\_trans}：分析最大 transition 违例
  \item \cmd{setup}：分析 setup 违例
  \item \cmd{hold}：分析 hold 违例
  \item \cmd{static\_noise}：分析静态噪声违例
  \item \cmd{only\_lib}：分析影响优化的库设置
\end{itemize}

此外，可指定：
\begin{itemize}
  \item 用 \opt{-fanout} 指定识别高扇出 net 的阈值。默认将扇出大于 30 的 net 视为高扇出 net。
  \item 用 \opt{-slack} 指定识别小的时序违例的松弛阈值。默认将小于 5~ps 的松弛违例视为小违例。
  \item 用 \opt{-drc\_nets} 指定用于最大 transition 分析的 net。默认分析所有 net。
  \item 用 \opt{-endpoints} 指定 setup/hold 分析的端点。默认分析所有端点。
  \item 用 \opt{-output} 指定输出文件名。工具会在你指定的文件名后添加 \texttt{.txt} 扩展名。默认生成名为 \texttt{analyze\_design\_violations.txt} 的文件。
\end{itemize}

% ============================================================
\section{分析布局}
\subsection*{Analyzing the Placement}

IC Compiler II 提供下列基于 GUI 与报告的能力，用于分析布局：
\begin{itemize}
  \item 报告利用率（Utilization）
  \item 报告布局 QoR
  \item 查询与更改布局状态
  \item 在 GUI 中分析布局
\end{itemize}

\subsection{报告利用率}
\subsubsection*{Reporting Utilization}

Block 的利用率按下列公式计算，其中 demand 为对象占用的总面积，capacity 为总可用面积：
\[
\text{Utilization} = \text{Demand} / \text{Capacity}
\]

报告利用率的步骤：
\begin{enumerate}
  \item （可选）用 \cmd{create\_utilization\_configuration} 创建利用率报告配置。\\
        下例为当前 block 创建名为 \cmd{config1} 的配置，用 core area 计算 capacity，并排除 hard 与 soft placement blockage：
\begin{lstlisting}
icc2_shell> create_utilization_configuration -scope block config1 \
   -capacity core_area -exclude {hard_blockages soft_blockages}
\end{lstlisting}
        要用 \cmd{remove\_utilization\_configurations} 移除利用率配置。
  \item 用 \cmd{report\_utilization} 报告利用率。\\
        用 \opt{-config} 与 \opt{-scope} 指定配置及其存储位置。\opt{-scope} 的有效值为 \cmd{tech}、\cmd{lib} 与 \cmd{block}。若与 \opt{-config} 联用时未指定 \opt{-scope}，工具依次在当前 block、设计库、逻辑库中搜索该配置。\\
        要用 \opt{-of\_objects} 报告特定对象的利用率，并指定一个或多个 block、voltage area、site array 或 bound。所有对象必须为同一类型。
\end{enumerate}

\subsection{报告布局 QoR}
\subsubsection*{Reporting the Placement QoR}

\cmd{report\_placement} 显示当前 block 的布局 QoR 信息。默认报告 block 中所有 net 的总半周长包围盒（half-perimeter boundary box）线长，以及任何布局违例，例如：
\begin{lstlisting}
icc2_shell> report_placement
...
...

  Wire length report (all)
  ==================
  wire length in design leon3mp: 12386597.894 microns.

  Physical hierarchy violations report
  ====================================
  Violations in design leon3mp:
     0 cells have placement violation.
\end{lstlisting}

要用 \opt{-hierarchical} 同时报告子 block 与当前 block。

要用 \opt{-wirelength} 及相应关键字，将线长计算限制为特定类型的 net：
\begin{itemize}
  \item \opt{-wirelength interface}：仅包含跨越物理边界的 net；
  \item \opt{-wirelength hard\_macro}：仅包含连接到 hard macro 的 net；
  \item \opt{-wirelength IO}：仅包含连接到 I/O pin 或 pad 的 net；
  \item \opt{-wirelength none}：抑制线长报告。
\end{itemize}

要用 \opt{-physical\_hierarchy\_violations} 及相应关键字，限制报告的布局违例类型：
\begin{itemize}
  \item \opt{-physical\_hierarchy\_violations internal}：仅报告放置在当前 block 布局区域之外的单元；
  \item \opt{-physical\_hierarchy\_violations external}：仅报告与子 block 重叠的单元；
  \item \opt{-physical\_hierarchy\_violations none}：抑制布局违例报告。
\end{itemize}

还可报告下列布局度量：
\begin{itemize}
  \item \textbf{Swimming pool area}\\
        Swimming pool 是由子 block、hard macro 与 placement blockage 围成的区域。此类区域可能使布局或布线难以成功，应覆盖 placement blockage 或予以消除。
        \begin{itemize}
          \item 用 \opt{-swimming\_pool\_area} 报告当前 block 中的总 swimming pool 面积；
          \item 用 \opt{-maximum\_swimming\_pool\_area} 报告面积小于特定阈值的 swimming pool 区域。
        \end{itemize}
  \item \textbf{Thin channel area}\\
        Thin channel 是当前 block 的子 block（无论是 hard macro 还是物理层次）之间的狭窄区域。此类通道可能使布局或布线难以成功，应覆盖 placement blockage 或加宽。
        \begin{itemize}
          \item 用 \opt{-thin\_channel\_area} 报告当前 block 中的总 thin channel 面积；
          \item 用 \opt{-maximum\_thin\_channel\_width} 或 \opt{-maximum\_thin\_channel\_height} 报告宽度或高度小于特定阈值的 thin channel。
        \end{itemize}
  \item \textbf{Hard macro 上方布线的线长}\\
        用 \opt{-hard\_macro\_route\_over} 报告当前 block 中在 hard macro 上方布线的 net 线长。该选项报告每个存在 over-macro 布线的 hard macro 的水平与垂直线长，以及 over-macro 布线的总线长。
  \item \textbf{Hard macro 重叠数量}\\
        用 \opt{-hard\_macro\_overlap} 报告当前 block 中重叠 hard macro 的总数。
\end{itemize}

\subsection{查询与更改布局状态}
\subsubsection*{Querying and Changing the Placement Status}

单元实例的 \cmd{status} 属性指示其当前布局状态。表~\ref{tab:place-status} 列出可能取值。

\begin{table}[htbp]
\centering
\caption{布局状态关键字}
\label{tab:place-status}
\begin{tabularx}{\textwidth}{@{}l X l@{}}
\toprule
\textbf{布局状态} & \textbf{说明} & \textbf{DEF 语法} \\
\midrule
\cmd{unplaced} & 单元实例尚未布局。
  注：要用 \cmd{reset\_placement} 将所有已布局单元恢复为 unplaced。该命令仅修改状态为 \cmd{placed} 的单元；不修改 \cmd{fixed} 或 \cmd{locked} 的单元。
  & \texttt{+UNPLACED} \\
\cmd{placed} & 单元实例已布局，但后续操作仍可移动。 & \texttt{+PLACED} \\
\cmd{fixed} & 单元实例已布局，后续操作可调整尺寸但不可移动。
  注：要防止这些单元被 sizing，须设置 \cmd{dont\_touch} 属性，见「优化期间保留单元与 Net」。
  & \texttt{+FIXED} \\
\cmd{legalize\_only} & 单元实例已布局，仅可由合法化器移动。 & N/A \\
\cmd{locked} & 单元实例已布局，不可手动或由工具移动。 & \texttt{+COVER} \\
\bottomrule
\end{tabularx}
\end{table}

要获取单元实例的布局状态，用 \cmd{get\_attribute} 查询其 \cmd{status} 属性。例如，获取名为 INST1 的单元实例：
\begin{lstlisting}
icc2_shell> get_attribute -objects [get_cells INST1] -name status
\end{lstlisting}

要更改单元实例的布局状态，使用 \cmd{set\_placement\_status}。例如，将 INST1 设为 \cmd{fixed}，使其在后续优化与布局中不被移动：
\begin{lstlisting}
icc2_shell> set_placement_status fixed [get_cells INST1]
\end{lstlisting}

\subsection{在 GUI 中分析布局}
\subsubsection*{Analyzing the Placement in the GUI}

IC Compiler II 图形用户界面（GUI）提供下列工具，用于分析与可视化布局后的结果质量：
\begin{itemize}
  \item \textbf{Global Route Congestion Map}：可视化布局在避免布线拥塞方面的质量。
  \item \textbf{Cell Density Map}：识别 block 中或矩形区域内的高单元密度区。
  \item \textbf{Pin Density Map}：识别 block 中或矩形区域内的高引脚密度区。
\end{itemize}

有关在 GUI 中使用这些分析工具的更多信息，见 \textit{IC Compiler II Graphical User Interface User Guide} 中的「Using Map and Visual Modes」。

% ============================================================
\section{分析时序}
\subsection*{Analyzing Timing}

IC Compiler II 会在需要时自动更新时序与延时估计结果。但可用 \cmd{update\_timing} 显式更新时序。

可用下列命令为 block 生成时序报告：
\begin{itemize}
  \item \cmd{report\_timing}：报告当前 block 的最差时序路径。
  \item \cmd{report\_qor}：显示当前 block 的 QoR 信息与统计。
  \item \cmd{report\_constraints}：报告当前 block 的逻辑 DRC 违例。
\end{itemize}

有关生成这些报告的更多信息，见 \textit{IC Compiler II Timing Analysis User Guide} 中的「Generating Reports」。

% ============================================================
\section{分析功耗}
\subsection*{Analyzing Power}

要计算并报告 block 的功耗，使用 \cmd{report\_power}。

\cmd{report\_power} 支持下列漏电功耗计算方法：
\begin{itemize}
  \item \textbf{Average}：基于所有状态的等权重概率。
  \item \textbf{Unconditional}：基于 \cmd{default\_cell\_leakage\_power} 库属性或 \cmd{cell\_leakage\_power} 库单元属性。
  \item \textbf{State-dependent}：基于单元各状态条件的加权和。计算取决于状态相关漏电库模型表征与静态概率。若状态相关漏电信息不可用，则功耗计算使用 unconditional 方法。
\end{itemize}

要为 \cmd{report\_power} 选择漏电功耗计算方法，将 \appopt{power.leakage\_mode} 设为 \cmd{average}、\cmd{unconditional} 或 \cmd{state}。要不报告漏电功耗，设为 \cmd{off}。默认为 \cmd{state}：
\begin{lstlisting}
icc2_shell> set_app_options -name power.leakage_mode -value average
\end{lstlisting}

要在功耗分析与报告期间使用复合电流源（CCS）接收端模型电容，将 \appopt{power.use\_ccs\_rcv\_cap} 设为 \cmd{true}。

下列示例给出默认功耗报告。

\begin{lstlisting}[caption={默认功耗报告}]
...
...
Mode: my_mode
Corner: bc_corner
Scenario: bc_scenario
Voltage: 1.160000
Temperature: 125.000000
Voltage unit: 1.00V
Temperature unit: 1.00C
Power unit: 1.00pW
Cell Leakage Power          = 5125495.58 uW
Power Group        Leakage Power          Total Power           (    % )
-----------------------------------------------------------------------------
---
io_pad                  0.00e+00             0.00e+00           ( 0.00%)
memory                  0.00e+00             0.00e+00           ( 0.00%)
black_box               0.00e+00             0.00e+00           ( 0.00%)
clock_network           3.32e+10             3.32e+01           ( 0.65%)
register                1.62e+12             1.62e+03           ( 31.67%)
sequential              1.16e+09             1.16e+00           ( 0.02%)
combinational           3.47e+12             3.47e+03           ( 67.66%)
-----------------------------------------------------------------------------
---
Total                   5.13e+12 pW          5.13e+03 mW
1
\end{lstlisting}

\subsection{为报告创建功耗组}
\subsubsection*{Creating Power Groups for Reporting}

运行 \cmd{report\_power} 且未指定单元列表时，工具报告默认功耗组的功耗信息。可用 \cmd{set\_power\_group} 创建由单元实例或库单元组成的功耗组。\cmd{report\_power} 会同时报告默认与用户指定功耗组的信息。

若单元属于多个功耗组，用于报告目的的归属按下列优先级确定：
\begin{enumerate}
  \item 用 \cmd{set\_power\_group} 为单元实例指定的功耗组；
  \item 用 \cmd{set\_power\_group} 为其库单元指定的功耗组；
  \item 该单元的默认功耗组。
\end{enumerate}

要用 \cmd{get\_power\_group} 获取特定单元的功耗组：
\begin{itemize}
  \item 用 \opt{-default} 获取其默认功耗组；
  \item 用 \opt{-user} 获取其可能所属的用户定义功耗组；
  \item 不带选项则获取用于报告目的的功耗组。
\end{itemize}

要用 \cmd{report\_power\_groups} 报告 block 中的所有功耗组。要用 \cmd{get\_power\_group\_objects} 获取特定功耗组中的单元集合。

要用 \cmd{reset\_power\_group} 移除功耗组，或从特定功耗组中移除单元。

\subsection{报告基于引脚的时钟网络功耗}
\subsubsection*{Reporting Pin-Based Clock Network Power}

\cmd{report\_power} 将时钟网络功耗作为名为 \cmd{clock\_network} 的默认功耗组的一部分进行报告。

默认情况下，报告的时钟网络功耗基于单元：若单元的某一输出在时钟网络上，则该单元的功耗（含其所有输入与输出的功耗）计入时钟网络。但可通过下列应用选项，将 \cmd{report\_power} 报告的时钟网络功耗限制为仅包含时钟网络上的引脚：
\begin{lstlisting}
icc2_shell> set_app_options -list {power.clock_network pin_based}
\end{lstlisting}

默认为 \cmd{cell\_based}。

\cmd{report\_clock\_qor -type power} 报告的功耗信息是基于引脚（pin-based）的。

% ============================================================
\section{对比 QoR 数据}
\subsection*{Comparing QoR Data}

可按下列步骤生成基于 Web 的报告，以用 QORsum 报告查看并对比 QoR 数据：
\begin{enumerate}
  \item （可选）用 \cmd{set\_qor\_data\_options} 配置后续步骤捕获与显示的 QoR 数据。
        \begin{itemize}
          \item 用 \opt{-leakage\_scenario} 与 \opt{-dynamic\_scenario} 指定设计中最关键的功耗 scenario。\\
                工具用你指定的功耗 scenario 生成 QORsum 报告中功耗 QoR 的高层摘要。若未指定，则对漏电与动态 scenario 均使用总功耗最高的 active 功耗 scenario。\\
                这些设置仅用于功耗 QoR 摘要；工具使用所有 active 功耗 scenario 的功耗信息来捕获并报告 QORsum 中的详细功耗信息。
          \item 用 \opt{-clock\_name} 与 \opt{-clock\_scenario} 指定最关键的时钟名与时钟 scenario。\\
                工具用你指定的时钟名与 scenario 生成时钟 QoR 的高层摘要。若未指定，工具识别最关键时钟并用于摘要。\\
                这些设置仅用于时钟 QoR 摘要；工具使用所有时钟生成详细时钟 QoR 信息。
          \item 用 \opt{-run\_name} 为 QORsum 报告中的该次运行指定标识名。\\
                默认工具用编号命名各次运行（如 Run1、Run2）。也可用更有意义的名称。在步骤 3 用 \cmd{compare\_qor\_data} 的 \opt{-run\_names} 也可指定运行名；若这样做，工具忽略 \cmd{set\_qor\_data\_options -run\_name} 指定的名称。
        \end{itemize}
        下例指定漏电功耗 scenario、动态功耗 scenario、时钟 scenario 以及用于对应摘要的时钟：
\begin{lstlisting}
icc2_shell> set_qor_data_options \
   -leakage_scenario S1 -dynamic_scenario S2 \
   -clock_scenario S3 -clock_name sys_clk
\end{lstlisting}

  \item 用 \cmd{write\_qor\_data} 收集报告用的 QoR 数据。该命令捕获与时序、功耗、运行时间及其他度量相关的 QoR 数据。\\
        可在希望捕获 QoR 数据的每个设计流程阶段多次运行 \cmd{write\_qor\_data}。用 \opt{-label} 指定值以标明捕获所处的流程阶段。\\
        默认 \cmd{write\_qor\_data} 捕获与已布局并优化设计相关的数据。可用 \opt{-report\_group} 及五个支持值之一（\cmd{unmapped}、\cmd{mapped}、\cmd{placed}、\cmd{cts}、\cmd{routed}）按流程状态调整捕获的报告集。更细粒度控制可用 \opt{-report\_list} 显式指定要生成的每个报告。\\
        \cmd{write\_qor\_data} 从多数常用工具报告（如 \cmd{report\_qor}、\cmd{report\_power}）向磁盘捕获标准 QoR 数据集。若有不属于命令集合的额外自定义数据要捕获，用 \cmd{capture\_qor\_data} 捕获完全自定义的 QoR 数据或 GUI 版图图像，以纳入基于 Web 的 QORsum 报告。这可用于补充 \cmd{write\_qor\_data} 捕获的标准 QoR 数据，或创建完全自定义的 QORsum 报告。相关命令包括：\cmd{define\_qor\_data\_panel}（在 QORsum 中创建新面板以查看 QoR 数据），以及 \cmd{set\_qor\_data\_metric\_properties}（自定义任意 QORsum 面板中任意度量列的属性，例如着色样式或对比数据时的着色阈值）。

  \item 用 \cmd{compare\_qor\_data} 生成 QORsum 报告。\\
        该命令取一次或多次 \cmd{write\_qor\_data} 运行捕获的数据，并创建基于 Web 的报告以查看与对比这些结果。\\
        必须用 \opt{-run\_locations} 指定上一步各次 \cmd{write\_qor\_data} 输出的位置；每个路径对应要对比的运行结果。\\
        可用 \opt{-run\_names} 为 QORsum 中的运行指定标识名；若这样做，工具忽略 \cmd{set\_qor\_data\_options -run\_name}。默认用编号命名（如 run1、run2）。\\
        可用 \opt{-output} 指定输出目录；默认写入名为 \texttt{compare\_qor\_data} 的目录。要用 \opt{-force} 覆盖已有输出数据；默认不覆盖。

  \item 用 \cmd{view\_qor\_data} 查看生成的 QORsum 报告。
\end{enumerate}

\begin{figure}[htbp]
\centering
\fbox{\parbox{0.85\textwidth}{\centering\vspace{1.2cm}
\textit{（原书 Figure 27：QORsum Report）}
\vspace{1.2cm}}}
\caption{QORsum 报告}
\label{fig:qorsum}
\end{figure}

有关探索对比数据的更多信息，见下列主题：
\begin{itemize}
  \item 设置基线运行（Baseline Run）
  \item 更改 QoR 显示样式
  \item 排序与过滤数据
  \item 探索详细对比数据
\end{itemize}

\subsection{设置基线运行}
\subsubsection*{Setting Your Baseline Run}

默认情况下，数据的第一行是基线运行，其他所有运行与之对比。基线运行的名称以金色高亮，标记为「golden」或基线结果。

非基线单元格的着色指示相对基线的差异方向与程度：
\begin{itemize}
  \item 红色表示相对基线变差；绿色表示相对基线改善；
  \item 较浅着色表示差异较小；较深着色表示差异较大。
\end{itemize}

要将光标悬停在列标题上，可查看各色阶对应的 delta 阈值。

要更改基线运行，单击 \textbf{Runs} 按钮，并从 \textbf{Baseline Run} 列选择新的基线运行。

\begin{seeAlsoBox}
更改 QoR 显示样式；排序与过滤数据；探索详细对比数据。
\end{seeAlsoBox}

\subsection{更改 QoR 显示样式}
\subsubsection*{Changing the QoR Display Style}

默认情况下，对比表显示所有运行的 QoR 值。例如，某次运行的 NVE 数为 2374。

在表中任意位置右键单击，可循环切换不同显示样式：
\begin{itemize}
  \item 相对基线的百分比 delta；
  \item 相对基线的绝对 delta（以斜体显示）。
\end{itemize}

例如，按百分比 delta，NVE 显示为比基线少 2.7\% 的失败端点；按绝对 delta，NVE 显示为比基线少 65 个失败端点。

\begin{seeAlsoBox}
排序与过滤数据；探索详细对比数据。
\end{seeAlsoBox}

\subsection{排序与过滤数据}
\subsubsection*{Sorting and Filtering the Data}

可对运行数据排序与过滤，以揭示结果中的模式，并确定要进一步探索的参数。

相关主题：
\begin{itemize}
  \item 排序数据
  \item 过滤度量（Metrics）
  \item 过滤运行（Runs）
  \item 分析示例
\end{itemize}

\begin{seeAlsoBox}
探索详细对比数据。
\end{seeAlsoBox}

\subsubsection{排序数据}
\subsubsection*{Sorting the Data}

单击任意列标题，按该度量排序。第一次单击升序，第二次单击降序。

例如，要使最差 TNS 位于表顶、最佳 TNS 位于表底，单击 TNS 列标题；再次单击则相反。

\subsubsection{过滤度量}
\subsubsection*{Filtering Metrics}

要控制表中显示哪些度量，单击 \textbf{Metrics} 按钮，并在 Metrics 对话框中相应勾选或取消勾选。

例如，要仅显示 setup 时序、网表面积与单元计数，按图~\ref{fig:metrics-dialog} 选择度量。

\begin{figure}[htbp]
\centering
\fbox{\parbox{0.85\textwidth}{\centering\vspace{1.0cm}
\textit{（原书 Figure 28：Metrics Dialog Box）}
\vspace{1.0cm}}}
\caption{Metrics 对话框}
\label{fig:metrics-dialog}
\end{figure}

\subsubsection{过滤运行}
\subsubsection*{Filtering Runs}

要控制表中显示哪些运行，单击 \textbf{Runs} 按钮，并从 \textbf{Visible Runs (Summary)} 列相应勾选或取消勾选探索运行。

例如，要仅显示前四次运行，按图~\ref{fig:choose-runs} 选择。

\begin{figure}[htbp]
\centering
\fbox{\parbox{0.85\textwidth}{\centering\vspace{1.0cm}
\textit{（原书 Figure 29：Choose Baseline and Visible Runs Dialog Box）}
\vspace{1.0cm}}}
\caption{选择基线与可见运行对话框}
\label{fig:choose-runs}
\end{figure}

\subsubsection{分析示例}
\subsubsection*{Example Analysis}

下例演示如何排序与过滤数据，以缩小要进一步探索的运行范围。

假设打开图~\ref{fig:cmp-report} 所示对比报告，并将 \cmd{util60\_aspect1\%1\_layerM9} 运行设为基线。

\begin{figure}[htbp]
\centering
\fbox{\parbox{0.85\textwidth}{\centering\vspace{1.0cm}
\textit{（原书 Figure 30：Comparison Report）}
\vspace{1.0cm}}}
\caption{对比报告}
\label{fig:cmp-report}
\end{figure}

可先查看 TNS，并将数据从最佳 TNS 排到最差 TNS。注意到最佳 TNS 运行使用 M9 顶层金属，最差使用 M7，说明限制金属层对时序影响显著。

随后可通过关闭 M7 运行的可见性，将分析限制为 M9 运行。在得到最佳 TNS 运行后，可按 GRCOverflow 列排序以将最差 overflow 置于顶部，从而对比拥塞。

注意到较高利用率运行的拥塞高于较低利用率运行。可按图~\ref{fig:lower-util} 关闭较高利用率运行的可见性，留下数量可控、更能满足时序与拥塞目标的探索运行子集。

\begin{figure}[htbp]
\centering
\fbox{\parbox{0.85\textwidth}{\centering\vspace{1.0cm}
\textit{（原书 Figure 31：Displaying Lower-Utilization Runs）}
\vspace{1.0cm}}}
\caption{显示较低利用率运行}
\label{fig:lower-util}
\end{figure}

\subsection{探索详细对比数据}
\subsubsection*{Exploring the Detailed Comparison Data}

启动 QORsum 报告时，应用打开 QOR Summary 表，汇总各次运行的高层时序、功耗与拥塞度量。该视图及其他摘要视图有助于排序与过滤数据，以定位要进一步探索的运行。

用摘要表识别最佳候选运行后，可用详细表获得更深洞察。摘要表给出所有运行的高层设计 QoR（如每次运行的整体设计时序 WNS、TNS、NVE），而详细表可显示设计中所有 path group 基于 path-group 的时序。详细表允许查看某一 path group 在所有运行上的时序，或一次查看最多六个运行的所有 path group 时序。所有详细表命名约定以「Details」结尾。例如，Path Type Details 表显示基于 path-type 的时序信息。

默认情况下，详细表从 Focused View 开始，如图~\ref{fig:focused-view}。

\begin{figure}[htbp]
\centering
\fbox{\parbox{0.85\textwidth}{\centering\vspace{1.0cm}
\textit{（原书 Figure 32：Focused View）}
\vspace{1.0cm}}}
\caption{Focused View}
\label{fig:focused-view}
\end{figure}

Focused View 一次显示一个 path group。要用顶部下拉菜单更改焦点。例如，可如图~\ref{fig:focused-dd} 指定 scenario 与 path group。

\begin{figure}[htbp]
\centering
\fbox{\parbox{0.85\textwidth}{\centering\vspace{1.0cm}
\textit{（原书 Figure 33：Focused View Drop-down Menu）}
\vspace{1.0cm}}}
\caption{Focused View 下拉菜单}
\label{fig:focused-dd}
\end{figure}

Focused View 聚焦于所有运行的某一特定属性。也可查看较少运行的全部属性。

例如，可同时查看所有 path group 以识别最差 path group。要切换到传统详细视图，单击顶部菜单的 Focused View 切换按钮，如图~\ref{fig:focused-toggle}。

\begin{figure}[htbp]
\centering
\fbox{\parbox{0.85\textwidth}{\centering\vspace{1.0cm}
\textit{（原书 Figure 34：Focused View Toggle Button）}
\vspace{1.0cm}}}
\caption{Focused View 切换按钮}
\label{fig:focused-toggle}
\end{figure}

要查看特定度量的详细对比数据，单击左侧面板中的某个详细视图。例如，要查看基于 path-group 的时序详情，选择 Path Group Details，如图~\ref{fig:path-group-details}。

\begin{figure}[htbp]
\centering
\fbox{\parbox{0.85\textwidth}{\centering\vspace{1.0cm}
\textit{（原书 Figure 35：Path Group Details）}
\vspace{1.0cm}}}
\caption{Path Group Details}
\label{fig:path-group-details}
\end{figure}

摘要表的每一行显示一次流程的结果，而详细表的每一行显示诸如 path group 名称等详细信息。各流程在每个度量下并排显示为列。例如，WNS 列组（以及每个度量的列组）显示标为 1、2、5、6、9、10 的相同六列。这些列代表各流程结果，用流程 ID 而非完整流程名，以避免列过宽。每个度量下的第一列是基线，其余五列为与之对比的测试流程。表上方有图例：金色框中的流程 ID 表示基线，蓝色框表示测试流程，后跟流程名。可单击 \textbf{Legend} 按钮展开或折叠图例。也可在单击 \textbf{Runs} 按钮打开的「Choose Baseline and Visible Runs」对话框中查看流程 ID。

详细表视图一次最多可显示六个运行（基线运行，加上在 Choose Baseline and Visible Runs 对话框中选择的前五个其他运行）。要更改显示的运行：
\begin{enumerate}
  \item 单击 \textbf{Runs} 按钮，打开 Choose Baseline and Visible Runs 对话框。
  \item 在 \textbf{Visible Runs (Summary)} 列中相应勾选或取消勾选运行。这会同步选中或取消选中 \textbf{Visible flows (Detailed)} 列中的那些运行。
\end{enumerate}

\subsubsection{过滤详细数据}
\subsubsection*{Filtering the Detailed Data}

详细视图提供过滤器以聚焦特定数据。某些视图有默认自动显示的过滤器。例如，下列详细视图具有默认的 scenario 与 path group 过滤器：Path Type Details、Path Group Details 与 Logic Level Details。

可通过下列方式修改默认过滤器：
\begin{itemize}
  \item 单击过滤器名称前的 X 符号以移除过滤器；
  \item 更改过滤器值；
  \item 单击过滤器上除 X 符号或值字段以外的任意位置，以启用或禁用过滤器。
\end{itemize}

也可对详细数据应用自定义过滤器。创建自定义过滤器：
\begin{enumerate}
  \item 单击 \textbf{Filter} 按钮，显示 Add Filter 字段。
  \item 通过定义过滤条件并选择要纳入结果的数据集来定义过滤器。\\
        选择列与比较类型，再指定比较值。
        \begin{itemize}
          \item 数值比较：选择 \texttt{=}、\texttt{!=}、\texttt{<}、\texttt{<=}、\texttt{>} 或 \texttt{>=}，并在 Value 字段中指定数值。
          \item 字符串比较：选择 \cmd{contains}，并指定字符串值。
          \item 正则表达式比较：选择 \cmd{regexp}，并指定正则表达式。
          \item 基于所选列可用值的过滤：选择 \cmd{enum}，Value 字段变为包含可用值的下拉菜单，再启用一个或多个显示的值。启用时，值前显示勾选标记。\\
                要启用或禁用某值：单击或用上下箭头导航到该值使其高亮，再按 Enter 切换状态。也可在 Value 字段中输入字符串以过滤下拉菜单中的可用值。单击 X 符号可取消已选值。
        \end{itemize}
  \item 单击绿色勾选标记应用过滤器。
\end{enumerate}

例如，要过滤 Path Group Details 视图以仅显示 \cmd{func@cworst} scenario 中的 path group，按相应界面定义过滤器。
